\corrigesChapitre{9}{Matrices et systèmes linéaires}
\begin{corrige}{1}{Dimensions}
\begin{questions}
\item $2\times4$, car les dimensions intérieures $3$ coïncident.
\item Non : les dimensions intérieures seraient $4$ et $2$.
\item Non : les tailles sont différentes.
\end{questions}
\end{corrige}
\begin{corrige}{2}{Addition et scalaire}
\begin{questions}
\item $\begin{pmatrix}2&4\\6&8\end{pmatrix}$.
\item $\begin{pmatrix}2&2\\3&5\end{pmatrix}$.
\item $I_2$, par convention sur les puissances.
\end{questions}
\end{corrige}
\begin{corrige}{3}{Produit ligne-colonne}
\begin{questions}
\item $AX=\binom{13}{15}$.
\item Une colonne de taille $2\times1$.
\item $2\times4+1\times5=13$, combinaison des coordonnées avec les coefficients de la première ligne.
\end{questions}
\end{corrige}
\begin{corrige}{4}{Inverser en ordre $2$}
\begin{questions}
\item $2-1=1$.
\item $\begin{pmatrix}1&-1\\-1&2\end{pmatrix}$.
\item $AA^{-1}=\begin{pmatrix}1&0\\0&1\end{pmatrix}$ par multiplication ligne-colonne.
\end{questions}
\end{corrige}
\begin{corrige}{5}{Non-commutativité}
\begin{questions}
\item $\begin{pmatrix}2&1\\4&3\end{pmatrix}$.
\item $\begin{pmatrix}3&4\\1&2\end{pmatrix}$.
\item Non : le développement général est $A^2+AB+BA+B^2$ ; il faut $AB=BA$ pour regrouper les termes croisés.
\end{questions}
\end{corrige}
\begin{corrige}{6}{Résoudre un système}
\begin{questions}
\item $A=\begin{pmatrix}2&1\\1&1\end{pmatrix}$, $X=\binom xy$, $Y=\binom53$.
\item $X=\begin{pmatrix}1&-1\\-1&2\end{pmatrix}\binom53=\binom21$.
\item $2(2)+1=5$ et $2+1=3$.
\end{questions}
\end{corrige}
\begin{corrige}{7}{Un ordre $3$}
\begin{questions}
\item En remontant : $x_3=2$, $x_2=3$, $x_1=0$.
\item $\begin{pmatrix}1&-1&1\\0&1&-1\\0&0&1\end{pmatrix}$.
\item Le coefficient $(1,3)$ vaut $1\times1+1\times(-1)+0\times1=0$, comme dans $I_3$.
\end{questions}
\end{corrige}
\begin{corrige}{8}{Rotation}
\begin{questions}
\item $R\binom21=\binom{-1}2$.
\item $R^2=-I_2$.
\item $R$ est la rotation d’angle $\pi/2$ ; $R^2$ est la symétrie centrale, ou rotation d’angle $\pi$.
\end{questions}
\end{corrige}
\begin{corrige}{9}{Matrice singulière}
\begin{questions}
\item $ad-bc=4-4=0$ ; la seconde ligne est le double de la première.
\item Les deux équations sont équivalentes à $x+2y=3$ : une infinité de solutions réelles.
\item Les équations se contredisent : la seconde devrait avoir second membre $6$, donc aucune solution.
\end{questions}
\end{corrige}
\begin{corrige}{10}{Ordre d’une composition}
\begin{questions}
\item $RS=\begin{pmatrix}0&1\\1&0\end{pmatrix}$ ; $SR=\begin{pmatrix}0&-1\\-1&0\end{pmatrix}$.
\item $R(SX)=(RS)X$ : la transformation appliquée en premier est à droite.
\item $RS$ donne $(1,1)$ ; $SR$ donne $(-1,-1)$, montrant la différence.
\end{questions}
\end{corrige}
\begin{corrige}{11}{Un inverse de produit}
\begin{questions}
\item Par associativité, $A(BB^{-1})A^{-1}=AA^{-1}=I$.
\item $B^{-1}(A^{-1}A)B=B^{-1}B=I$.
\item $(AB)^{-1}=B^{-1}A^{-1}$. L’ordre des facteurs est inversé ; $A^{-1}B^{-1}$ ne convient pas en général.
\end{questions}
\end{corrige}
\begin{corrige}{12}{Translation et matrice}
\begin{questions}
\item Chaque coefficient du produit est une somme de produits par zéro, donc nul.
\item Elle déplace l’origine, contrairement à toute application $X\mapsto AX$.
\item $X\mapsto I_2X+\binom3{-1}$.
\end{questions}
\end{corrige}
